\documentclass{article}
\usepackage[margin=1in]{geometry}
\title{Compact Feature Representations for Synthetic Document Classification}
\author{Anonymous}
\date{}
\begin{document}
\maketitle

\begin{abstract}
We evaluate a compact feature representation for classifying deterministic synthetic documents into two categories. A logistic-regression classifier trained on 16-dimensional features obtains 0.90 accuracy on a held-out test split.
\end{abstract}

\section{Introduction}
We study a controlled document-classification setting in which all inputs are synthetic. The purpose is to validate a reproducible paper-writing task contract rather than make a claim about a real-world corpus.

\section{Related Work}
Reproducible agent evaluation requires explicit task inputs and stable output artifacts \cite{harbor2024}.

\section{Method}
We generate 1,000 deterministic synthetic documents. Each document is represented by a 16-dimensional feature vector, and a logistic-regression classifier predicts one of two categories.

\section{Experiments}
We reserve 20 percent of the documents for testing and train on the remaining 80 percent. Accuracy on the held-out partition is the sole reported metric.

\section{Results}
The classifier obtains 0.90 held-out accuracy. This number is illustrative and applies only to the synthetic setup.

\section{Conclusion}
The result demonstrates that the supplied materials support a complete, compilable paper. The task does not establish performance on real-world data.

\bibliographystyle{plain}
\bibliography{references}
\end{document}
